\documentclass[11pt,a4paper]{article}
\usepackage[utf8]{inputenc}
\usepackage[T1]{fontenc}
\usepackage[french]{babel}
\usepackage{geometry}
\usepackage{xcolor}
\usepackage{hyperref}
\usepackage{titlesec}

\geometry{margin=1in}
\definecolor{titleblue}{RGB}{0,51,102}

\hypersetup{
    colorlinks=true,
    linkcolor=titleblue,
    urlcolor=titleblue,
    pdftitle={Lettre de Motivation - Ismail AISSAOUI},
}

\begin{document}

\noindent \textbf{Ismail AISSAOUI} \\
Ingénieur d'État en Intelligence Artificielle (ESI-SBA) \\
Email: ismail.aissaoui.pro@gmail.com \\
Téléphone: +213 660 70 77 96 \\
LinkedIn: https://ismail-aissaoui.vercel.app

\bigskip

\begin{flushright}
    Au Comité de Sélection du Doctorat \\
    \textbf{Programme EDT - IMT Atlantique / IRISA} \\
    À l'attention de : Prof. Antoine Beugnard et Prof. Jean-Marc Jézéquel \\
    \medskip
    1er mai 2026
\end{flushright}

\bigskip

\noindent \textbf{Objet : Candidature au doctorat : Fédération de données et de modèles pour définir des jumeaux numériques (Rang 6 - PC2)}

\bigskip

Monsieur le Professeur,

Ingénieur d'État en Intelligence Artificielle diplômé de l'ESI-SBA, je vous adresse avec enthousiasme ma candidature pour le poste de doctorat intitulé « Fédération de données et de modèles pour définir des jumeaux numériques » au sein du programme national Engineering Digital Twin (EDT). Cette candidature est motivée par l'adéquation de mon profil avec vos travaux sur la fédération d'informations et l'ingénierie des modèles.

Actuellement en fin de cycle ingénieur chez Sonatrach, je travaille sur le développement de **Jumeaux Numériques informés par la physique** qui intègrent des flux de capteurs hétérogènes, des modèles CAO 3D et des équations physiques. Cette expérience m'a confronté au défi majeur de l'hétérogénéité des données et de leur synchronisation temporelle. Je suis particulièrement motivé par votre proposition d'utiliser la fédération d'informations et les ontologies pour gérer l'espace multidimensionnel des variantes et du temps dans les référentiels de jumeaux numériques.

Votre expertise mondiale en Ingénierie Dirigée par les Modèles (MDE) et en architecture logicielle représente pour moi un environnement de recherche idéal. Je suis convaincu que ma maîtrise de l'extraction de connaissances (RAG) et ma compréhension des problématiques de cycle de vie des données industrielles me permettront de contribuer de manière significative à vos objectifs de définition de jumeaux numériques par assemblage.

Je suis un chercheur autonome et motivé, capable de faire le pont entre les architectures logicielles complexes et les besoins opérationnels de la **plateforme Artemis**.

Je vous remercie de l'attention que vous porterez à ma candidature et reste à votre entière disposition pour un entretien.

Dans l'attente de votre réponse, je vous prie d'agréer, Monsieur le Professeur, l'expression de mes salutations distinguées.

\bigskip

\noindent Ismail AISSAOUI

\end{document}
